\begin{definition}[Bond-two vertical tensor and Bell matrix]%
    \label{def:mpdo_bond_two_singleton_deformation}
    \lean{MPOTensor.BondTwoSingletonGramBoundary.singletonTensor}
    \lean{MPOTensor.BondTwoSingletonGramBoundary.terminalJ}
    \lean{MPOTensor.BondTwoSingletonGramBoundary.bellVector}
    \lean{MPOTensor.BondTwoSingletonGramBoundary.bellProjector}
    \lean{MPOTensor.BondTwoSingletonGramBoundary.gaugeMatrix}
    \lean{MPOTensor.BondTwoSingletonGramBoundary.gauge}
    \lean{MPOTensor.BondTwoSingletonGramBoundary.gaugedSingletonTensor}
    \lean{MPOTensor.BondTwoSingletonGramBoundary.deformedTerminal}
    \lean{MPOTensor.BondTwoSingletonGramBoundary.deformedCompanionBell}
    \leanok
    \uses{def:mpo_physical_trace_transfer}
    Let $I=\{0,1\}$ and let $E_{a,b}$ be the matrix units of
    $\MN{2}$. Define a tensor $A$ by
    \begin{align}
        A^{((a,b),(c,d))}
         & =\begin{cases}
                E_{a,b}, & (a,b)=(c,d),   \\
                0,       & (a,b)\ne(c,d).
            \end{cases}
        \notag
    \end{align}
    Its physical trace is $J=\begin{pmatrix}1&1\\1&1\end{pmatrix}$.
    Define $\ket{\Phi}=\ket{0,0}+\ket{1,1}$ and
    $P=\ketbra{\Phi}{\Phi}$. Thus $P$ is an unnormalized rank-one positive
    operator, not an orthogonal projection.
    The tensor $A$ is the vertical tensor of the physical model in
    Definition~\ref{def:mpdo_bond_two_singleton_base_mpo}.
    Theorem~\ref{thm:mpdo_bond_two_singleton_bell_identity} identifies $P$
    with its two-site operator after reindexing.
    For $X=\frac12\begin{pmatrix}3&1\\1&3\end{pmatrix}$, set
    $B^v=XA^vX^{-1}$ and define
    \begin{align}
        J_X & =XJX^{-1},
        \notag                                      \\
        P_X & =(X\otimes X)P(X^{-1}\otimes X^{-1}).
        \notag
    \end{align}
\end{definition}

\begin{lemma}[Positive physical trace and non-Hermitian Bell similarity]%
    \label{lem:mpdo_bond_two_singleton_nonunitary_deformation}
    \lean{MPOTensor.BondTwoSingletonGramBoundary.singletonTensor_isInjective}
    \lean{MPOTensor.BondTwoSingletonGramBoundary.singletonTensor_isNormal}
    \lean{MPOTensor.BondTwoSingletonGramBoundary.physTraceTransfer_singletonTensor}
    \lean{MPOTensor.BondTwoSingletonGramBoundary.terminalJ_posSemidef}
    \lean{MPOTensor.BondTwoSingletonGramBoundary.bellProjector_posSemidef}
    \lean{MPOTensor.BondTwoSingletonGramBoundary.gauge_gram_ne_one}
    \lean{MPOTensor.BondTwoSingletonGramBoundary.%
        singletonTensor_gaugeEquiv_gaugedSingletonTensor}
    \lean{MPOTensor.BondTwoSingletonGramBoundary.%
        singletonTensor_sameMPV₂Pos_gaugedSingletonTensor}
    \lean{MPOTensor.BondTwoSingletonGramBoundary.gauge_commutes_terminalJ}
    \lean{MPOTensor.BondTwoSingletonGramBoundary.%
        physTraceTransfer_gaugedSingletonTensor}
    \lean{MPOTensor.BondTwoSingletonGramBoundary.deformedTerminal_eq_terminalJ}
    \lean{MPOTensor.BondTwoSingletonGramBoundary.deformedTerminal_posSemidef}
    \lean{MPOTensor.BondTwoSingletonGramBoundary.deformedCompanionBell_not_isHermitian}
    \leanok
    \uses{def:mpdo_bond_two_singleton_deformation}
    The tensor $A$ is injective and proportional to an NT. The matrix $X$
    is invertible but not unitary, and $A$ and $B$ generate the same matrix
    product vectors at every positive length. Their physical traces satisfy
    \begin{align}
        \sum_{u\in I\times I} A^{(u,u)}
            & =\sum_{u\in I\times I} B^{(u,u)}=J\succeq0,
        \notag                                            \\
        J_X & =J.
        \notag
    \end{align}
    Nevertheless, $P$ is positive semidefinite while $P_X$ is not Hermitian.
    Theorem~\ref{thm:mpdo_bond_two_singleton_bell_identity} identifies $P$
    with the reindexed two-site operator of the physical model. Equality
    of the vertical matrix product vectors does not imply that the
    corresponding physical deformation generates MPDOs.
\end{lemma}

\begin{proof}\leanok
    \uses{def:mpdo_bond_two_singleton_deformation}
    The four diagonal physical letters of $A$ are the matrix units, so $A$ is
    injective. Similarity preserves every cyclic contraction. Direct
    multiplication gives $XJ=JX$, hence $J_X=J\succeq0$, whereas
    $X^\dagger X\ne\Id$. In the Bell calculation, the $(00,01)$ and
    $(01,00)$ entries of the real matrix $P_X$ are unequal, so $P_X$ is not
    Hermitian.
\end{proof}

\begin{theorem}[Physical-trace and Bell Hermiticity force a scalar Gram matrix]%
    \label{thm:mpdo_bond_two_singleton_scalar_gram}
    \lean{MPOTensor.BondTwoSingletonGramBoundary.%
        posDef_eq_pos_smul_one_of_commutes_terminalJ_of_bellCross}
    \lean{MPOTensor.BondTwoSingletonGramBoundary.%
        gaugeGram_commutes_terminalJ_of_terminal_isHermitian}
    \lean{MPOTensor.BondTwoSingletonGramBoundary.%
        gaugeGram_bellCross_of_companionBell_isHermitian}
    \lean{MPOTensor.BondTwoSingletonGramBoundary.%
        gaugeGram_eq_pos_smul_one_of_terminal_companionBell_isHermitian}
    \leanok
    \uses{def:mpdo_bond_two_singleton_deformation}
    Let $Y\in\GL(2,\C)$. If $YJY^{-1}$ and
    $(Y\otimes Y)P(Y^{-1}\otimes Y^{-1})$ are Hermitian, then there is a
    real number $\omega>0$ such that $Y^\dagger Y=\omega\Id$.
    The comparisons in
    \cite[Appendix~C.4, lines~2048--2057]{Cirac2016MPDO_arXiv} and
    \cite[Proposition~4.13]{Cirac2016MPDO_arXiv} motivate this
    finite-dimensional criterion; neither source states this model-specific
    theorem. Theorem~\ref{thm:mpdo_bond_two_singleton_bell_identity} identifies
    $P$ with the two-site operator of the undeformed physical model. The
    Hermiticity assumptions here do not assert that a corresponding
    deformation generates MPDOs at every length.
\end{theorem}

\begin{proof}\leanok
    \uses{def:mpdo_bond_two_singleton_deformation}
    Put $G=Y^\dagger Y$, which is positive definite. Hermiticity of
    $YJY^{-1}$ gives $GJ=JG$. Hermiticity of the Bell similarity gives
    $G_{00}G_{01}+G_{10}G_{11}=0$. The commutator identity implies
    $G_{01}=G_{10}$ and $G_{00}=G_{11}$. Since $G_{00}>0$, the remaining
    equation forces $G_{01}=G_{10}=0$, and therefore $G=G_{00}\Id$.
\end{proof}

\subsection{A physical MPDO with a single vertical BNT element}

\begin{definition}[Physical tensor with a bond-two vertical tensor]%
    \label{def:mpdo_bond_two_singleton_base_mpo}
    \lean{MPOTensor.BondTwoSingletonBaseModel.baseMPO}
    \leanok
    \uses{def:mpo_tensor, def:mpdo_bond_two_singleton_deformation}
    Let $I=\{0,1\}$ and identify $I\times I$ with four bond coordinates.
    Define the bond-four MPO tensor $M$ by
    \begin{align}
        M^{i,j} & =E_{(i,j),(i,j)}.
        \notag
    \end{align}
    This model is motivated by
    \cite[Appendix~C.4, lines~2048--2057]{Cirac2016MPDO_arXiv}; it is not a
    tensor stated there.
\end{definition}

\begin{theorem}[Vertical tensor identity]%
    \label{thm:mpdo_bond_two_singleton_vertical_identity}
    \lean{MPOTensor.BondTwoSingletonBaseModel.verticalTensor_baseMPO}
    \leanok
    \uses{def:mpdo_bond_two_singleton_base_mpo, def:mpdo_bond_two_singleton_deformation}
    The vertical tensor of $M$ is exactly the matrix-unit tensor $A$ from
    Definition~\ref{def:mpdo_bond_two_singleton_deformation}.
\end{theorem}
\begin{proof}\leanok
    A doubled bond coordinate $(a,b)$ selects $E_{a,b}$ when the two physical
    pair labels agree, and gives zero otherwise.
\end{proof}

\begin{theorem}[Positive-length GHZ formula in the physical direction]%
    \label{thm:mpdo_bond_two_singleton_ambient_ghz}
    \lean{MPOTensor.BondTwoSingletonBaseModel.constantPairWeight}
    \lean{MPOTensor.BondTwoSingletonBaseModel.ghzAmplitude}
    \lean{MPOTensor.BondTwoSingletonBaseModel.unnormalizedGHZRankOne}
    \lean{MPOTensor.BondTwoSingletonBaseModel.mpo_baseMPO_zero}
    \lean{MPOTensor.BondTwoSingletonBaseModel.unnormalizedGHZRankOne_zero}
    \lean{MPOTensor.BondTwoSingletonBaseModel.%
        mpo_baseMPO_zero_ne_unnormalizedGHZRankOne_zero}
    \lean{MPOTensor.BondTwoSingletonBaseModel.%
        baseMPO_eq_cyclicEdgeWeightTensor}
    \lean{MPOTensor.BondTwoSingletonBaseModel.%
        mpo_baseMPO_eq_unnormalizedGHZRankOne}
    \lean{MPOTensor.BondTwoSingletonBaseModel.baseMPO_isMPDO}
    \leanok
    \uses{def:mpdo_bond_two_singleton_base_mpo}
    For every $N\geq1$, let
    $\ket{\mathrm{GHZ}_N}=\ket{0}^{\otimes N}+\ket{1}^{\otimes N}$.
    Then
    \begin{align}
        \rho^{(N)}(M)
         & =\ketbra{\mathrm{GHZ}_N}{\mathrm{GHZ}_N}.
        \notag
    \end{align}
    Hence $M$ generates MPDOs. These are unnormalized rank-one positive
    operators, not orthogonal projections. At $N=0$, the closed MPO
    contraction instead gives $4\Id$, whereas the empty-word outer product
    is defined to be $\Id$. Thus the positive-length restriction is necessary.
\end{theorem}
\begin{proof}\leanok
    The cyclic equality weights force the ket and bra words to be independently
    constant. The result is an outer product and is therefore positive
    semidefinite.
\end{proof}

\begin{theorem}[Two-site Bell identity under reindexing]%
    \label{thm:mpdo_bond_two_singleton_bell_identity}
    \lean{MPOTensor.BondTwoSingletonBaseModel.%
        twoSite_baseMPO_eq_bellProjector}
    \leanok
    \uses{thm:mpdo_bond_two_singleton_ambient_ghz, def:mpdo_bond_two_singleton_deformation}
    Let $\varepsilon$ identify a binary word of length two with its ordered
    pair of letters, and put $\ket{\Phi}=\ket{0,0}+\ket{1,1}$. Then
    \begin{align}
        \reindex_{\varepsilon,\varepsilon}
        (\rho^{(2)}(M))
         & =\ketbra{\Phi}{\Phi}.
        \notag
    \end{align}
    Thus $P$ is exactly the two-site physical operator in ordered-pair
    coordinates.
\end{theorem}
\begin{proof}\leanok
    Specialize Theorem~\ref{thm:mpdo_bond_two_singleton_ambient_ghz} to
    $N=2$ and evaluate the coordinate identification.
\end{proof}

\begin{definition}[Physical similarity of the MPO tensor]%
    \label{def:mpdo_bond_two_singleton_physical_deformation}
    \lean{MPOTensor.BondTwoSingletonBaseModel.gaugeDeformedBaseMPO}
    \leanok
    \uses{def:mpdo_bond_two_singleton_base_mpo}
    For $X\in\GL(2,\C)$, define $M_X$ by applying $X$ to the ket index and
    $X^{-1}$ to the bra index of $M$:
    \begin{align}
        M_X^{i,j}
         & =\sum_{a,b\in I}X_{i,a}(X^{-1})_{b,j}M^{a,b}.
        \notag
    \end{align}
    This family is motivated by the sector comparison in
    \cite[Appendix~C.4, lines~2048--2057]{Cirac2016MPDO_arXiv} and the
    physical-sector argument of
    \cite[Proposition~4.13, lines~1898--1921]{Cirac2016MPDO_arXiv}.
    Neither passage states this deformation or the results below.
\end{definition}

\begin{theorem}[Vertical tensor under physical similarity]%
    \label{thm:mpdo_bond_two_singleton_physical_vertical}
    \lean{MPOTensor.BondTwoSingletonBaseModel.%
        verticalTensor_gaugeDeformedBaseMPO}
    \leanok
    \uses{def:mpdo_bond_two_singleton_physical_deformation, thm:mpdo_bond_two_singleton_vertical_identity}
    For every vertical letter $v$, the vertical tensor of $M_X$ satisfies
    $\widetilde M_X^v=XA^vX^{-1}$.
\end{theorem}
\begin{proof}\leanok
    \uses{def:mpdo_vertical_one_site_actions, thm:mpdo_bond_two_singleton_vertical_identity}
    Left and right multiplication on the physical indices become matrix
    multiplication on the vertically viewed tensor. Now use
    Theorem~\ref{thm:mpdo_bond_two_singleton_vertical_identity}.
\end{proof}

\begin{theorem}[One-site physical-similarity identity]%
    \label{thm:mpdo_bond_two_singleton_physical_one_site}
    \lean{MPOTensor.BondTwoSingletonBaseModel.%
        oneSite_gaugeDeformedBaseMPO_eq_terminalJ}
    \leanok
    \uses{def:mpdo_bond_two_singleton_physical_deformation, def:mpdo_bond_two_singleton_deformation}
    After identifying one-site words with their single letters in $I$, the
    one-site operator of $M_X$ is
    \begin{align}
        \reindex(\rho^{(1)}(M_X)) & =XJX^{-1}.
        \notag
    \end{align}
\end{theorem}
\begin{proof}\leanok
    Expand the one-site contraction. The four diagonal physical letters sum to
    $J$, and the two physical multiplications give conjugation by $X$.
\end{proof}

\begin{theorem}[Two-site physical-similarity identity]%
    \label{thm:mpdo_bond_two_singleton_physical_two_site}
    \lean{MPOTensor.BondTwoSingletonBaseModel.%
        twoSite_gaugeDeformedBaseMPO_eq_bellProjector}
    \leanok
    \uses{def:mpdo_bond_two_singleton_physical_deformation, thm:mpdo_bond_two_singleton_bell_identity}
    Under the ordered-pair identification $\varepsilon$ from
    Theorem~\ref{thm:mpdo_bond_two_singleton_bell_identity},
    \begin{align}
        \reindex_{\varepsilon,\varepsilon}
        (\rho^{(2)}(M_X))
         & =(X\otimes X)P(X^{-1}\otimes X^{-1}).
        \notag
    \end{align}
\end{theorem}
\begin{proof}\leanok
    Evaluate the two-site contraction in ordered-pair coordinates. Each ket
    index contributes $X$, each bra index contributes $X^{-1}$, and the
    undeformed contraction is $P$.
\end{proof}

\begin{theorem}[MPDO positivity forces a scalar physical Gram matrix]%
    \label{thm:mpdo_bond_two_singleton_physical_scalar_gram}
    \lean{MPOTensor.BondTwoSingletonBaseModel.%
        gaugeGram_eq_pos_smul_one_of_gaugeDeformedBaseMPO_isMPDO}
    \leanok
    \uses{thm:mpdo_bond_two_singleton_physical_one_site, thm:mpdo_bond_two_singleton_physical_two_site, thm:mpdo_bond_two_singleton_scalar_gram}
    If $M_X$ generates MPDOs, then there is a real number $\omega>0$ such
    that $X^\dagger X=\omega\Id$. Thus a physical similarity with
    nonscalar $X^\dagger X$ does not preserve MPDO positivity in this family.
\end{theorem}
\begin{proof}\leanok
    \uses{thm:mpdo_bond_two_singleton_physical_one_site, thm:mpdo_bond_two_singleton_physical_two_site, thm:mpdo_bond_two_singleton_scalar_gram}
    Positivity of the one- and two-site operators implies Hermiticity. Apply
    the two exact reindexing identities and then
    Theorem~\ref{thm:mpdo_bond_two_singleton_scalar_gram}.
\end{proof}

\begin{lemma}[The concrete physical Gram matrix is nonscalar]%
    \label{lem:mpdo_bond_two_singleton_physical_nonscalar_gram}
    \lean{MPOTensor.BondTwoSingletonBaseModel.gauge_gram_ne_smul_one}
    \leanok
    \uses{def:mpdo_bond_two_singleton_deformation}
    For $X=\frac12\begin{pmatrix}3&1\\1&3\end{pmatrix}$ and every
    $\omega\in\R$, we have $X^\dagger X\ne\omega\Id$.
\end{lemma}
\begin{proof}\leanok
    The $(0,1)$ entry of $X^\dagger X$ equals $3/2$, whereas every scalar
    matrix has vanishing $(0,1)$ entry.
\end{proof}

\begin{theorem}[The concrete physical similarity does not generate MPDOs]%
    \label{thm:mpdo_bond_two_singleton_physical_not_mpdo}
    \lean{MPOTensor.BondTwoSingletonBaseModel.%
        gaugeDeformedBaseMPO_gauge_not_isMPDO}
    \leanok
    \uses{thm:mpdo_bond_two_singleton_physical_scalar_gram, lem:mpdo_bond_two_singleton_physical_nonscalar_gram}
    For $X=\frac12\begin{pmatrix}3&1\\1&3\end{pmatrix}$, the physical
    deformation $M_X$ does not generate MPDOs.
\end{theorem}
\begin{proof}\leanok
    \uses{thm:mpdo_bond_two_singleton_physical_scalar_gram, lem:mpdo_bond_two_singleton_physical_nonscalar_gram}
    Otherwise Theorem~\ref{thm:mpdo_bond_two_singleton_physical_scalar_gram}
    would make $X^\dagger X$ scalar, contrary to
    Lemma~\ref{lem:mpdo_bond_two_singleton_physical_nonscalar_gram}.
\end{proof}

\begin{theorem}[Canonical form in the original physical direction]%
    \label{thm:mpdo_bond_two_singleton_cpsv_form}
    \lean{MPOTensor.BondTwoSingletonBaseModel.baseSymbolTensor}
    \lean{MPOTensor.BondTwoSingletonBaseModel.baseSymbolTensor_isNormalTensor}
    \lean{MPOTensor.BondTwoSingletonBaseModel.baseSymbolBlocks}
    \lean{MPOTensor.BondTwoSingletonBaseModel.%
        baseMPO_toMPSTensor_eq_symbolBlocks}
    \lean{MPOTensor.BondTwoSingletonBaseModel.%
        baseMPO_toMPSTensor_isCPSVCanonicalForm}
    \leanok
    \uses{def:mpdo_bond_two_singleton_base_mpo}
    As an MPS with four doubled physical symbols, $M$ is the direct sum of
    four bond-one NTs, one supported on each symbol, all with weight one.
    Hence $M$ itself is in CF, without a change of basis or blocking.
\end{theorem}
\begin{proof}\leanok
    Each block has identity transfer map, and their direct sum is the diagonal
    matrix-unit tensor $M$.
\end{proof}

\begin{theorem}[Cyclic-edge projection in the vertical direction]%
    \label{thm:mpdo_bond_two_singleton_retained_projection}
    \lean{MPOTensor.BondTwoSingletonBaseModel.edgeSource}
    \lean{MPOTensor.BondTwoSingletonBaseModel.edgeTarget}
    \lean{MPOTensor.BondTwoSingletonBaseModel.retainedCyclicIndicator}
    \lean{MPOTensor.BondTwoSingletonBaseModel.%
        mpo_verticalBNTMPO_singletonTensor_zero}
    \lean{MPOTensor.BondTwoSingletonBaseModel.%
        retainedCyclicIndicator_zero_diagonal}
    \lean{MPOTensor.BondTwoSingletonBaseModel.%
        mpo_verticalBNTMPO_singletonTensor_zero_ne_diagonal}
    \lean{MPOTensor.BondTwoSingletonBaseModel.%
        mpo_verticalBNTMPO_singletonTensor_eq_diagonal}
    \lean{MPOTensor.BondTwoSingletonBaseModel.%
        mpo_verticalBNTMPO_singletonTensor_isHermitian_and_idempotent}
    \lean{MPOTensor.BondTwoSingletonBaseModel.%
        mpo_verticalBNTMPO_singletonTensor_isStarProjection}
    \lean{MPOTensor.BondTwoSingletonBaseModel.%
        raw_singleton_coefficient_one_algebra}
    \leanok
    \uses{thm:mpdo_bond_two_singleton_vertical_identity}
    Encode the operator indices of the vertical tensor $A$ as directed
    edges on $I$. For $L\geq1$, let $q_L(w)$ indicate that consecutive
    edges in $w$ meet cyclically. The closed operator in this direction is
    \begin{align}
        O_L(A)         & =\diag(q_L),
        \notag                             \\
        O_L(A)^\dagger & =O_L(A)=O_L(A)^2.
        \notag
    \end{align}
    Thus $O_L(A)$ is an orthogonal projection and satisfies an algebra law
    with one label and structure coefficient one. It is distinct from the
    unnormalized GHZ rank-one operator in the original physical direction
    in Theorem~\ref{thm:mpdo_bond_two_singleton_ambient_ghz}.
    At $L=0$, the vertical contraction is $O_0(A)=2\Id$, whereas the empty
    cyclic indicator gives $\Id$. In particular, this contraction is not
    idempotent at length zero, so neither the projection identity nor the
    coefficient-one algebra law extends to $L=0$.
\end{theorem}
\begin{proof}\leanok
    The cyclic indicator takes only the values zero and one, so its diagonal
    matrix is Hermitian and idempotent.
\end{proof}

\begin{theorem}[Normalized one-element vertical BNT and its algebra]%
    \label{thm:mpdo_bond_two_singleton_normalized_algebra}
    \lean{MPOTensor.BondTwoSingletonBaseModel.singletonScale}
    \lean{MPOTensor.BondTwoSingletonBaseModel.normalizedSingletonTensor}
    \lean{MPOTensor.BondTwoSingletonBaseModel.%
        normalizedSingletonTensor_isLeftCanonical}
    \lean{MPOTensor.BondTwoSingletonBaseModel.%
        normalizedSingletonTensor_isNormalTensor}
    \lean{MPOTensor.BondTwoSingletonBaseModel.%
        mpo_verticalBNTMPO_normalizedSingletonTensor}
    \lean{MPOTensor.BondTwoSingletonBaseModel.%
        normalized_singleton_geometric_algebra}
    \lean{MPOTensor.BondTwoSingletonBaseModel.normalizedSingletonChi}
    \lean{MPOTensor.BondTwoSingletonBaseModel.normalizedSingletonCoeffs}
    \lean{MPOTensor.BondTwoSingletonBaseModel.%
        normalizedSingletonCoeffs_coeff}
    \lean{MPOTensor.BondTwoSingletonBaseModel.%
        normalizedSingleton_isCPSVBasis}
    \lean{MPOTensor.BondTwoSingletonBaseModel.%
        normalizedSingleton_verticalAssembledTensor}
    \lean{MPOTensor.BondTwoSingletonBaseModel.%
        normalizedSingletonAlgebraClause}
    \leanok
    \uses{thm:mpdo_bond_two_singleton_retained_projection}
    Put $s=2^{-1/2}$ and $\widehat A=sA$. The tensor $\widehat A$ is
    left-canonical and is an NT with transfer spectral radius one.
    A vertical BNT consists of the single tensor $M_0=\widehat A$.
    With one copy, take the positive diagonal matrices
    $\mu_0=(s^{-1})$ and $\chi_{0,0,0}=(s)$. The structure coefficient is
    $c_{0,0,0}^{(L)}=\tr(\chi_{0,0,0}^L)=s^L$. For $L\geq1$,
    \begin{align}
        O_L(M_0)   & =s^L O_L(A),
        \notag                      \\
        O_L(M_0)^2 & =s^L O_L(M_0).
        \notag
    \end{align}
    The weighted tensor $\mu_0\otimes M_0$ reconstructs $A$ exactly.
    For $m_0=\tr(\mu_0)=s^{-1}$, the identity
    $m_0=c_{0,0,0}^{(1)}m_0^2$ says that the vector $m$ is an idempotent
    for the multiplication induced by $c^{(1)}$.
    At $L=0$, however, $O_0(M_0)=2\Id$, so
    $O_0(M_0)^2=4\Id\ne2\Id$; the algebra identity therefore requires
    positive length.
\end{theorem}
\begin{proof}\leanok
    The identity $s^2=1/2$ gives left-canonical normalization. Scaling every
    local tensor contributes $s^L$ to a length-$L$ contraction, and $s^{-1}$
    cancels the one-site scaling in the weighted reconstruction.
\end{proof}

\begin{theorem}[Coisometry identifying the vertical sector coordinates]%
    \label{thm:mpdo_bond_two_singleton_coisometry}
    \lean{MPOTensor.BondTwoSingletonBaseModel.singletonSectorCoordinateEquiv}
    \lean{MPOTensor.BondTwoSingletonBaseModel.singletonRetainedCoordinateEquiv}
    \lean{MPOTensor.BondTwoSingletonBaseModel.singletonVerticalCoisometry}
    \lean{MPOTensor.BondTwoSingletonBaseModel.%
        singletonVerticalCoisometry_coisometry}
    \lean{MPOTensor.BondTwoSingletonBaseModel.%
        singletonVerticalCoisometry_conj}
    \lean{MPOTensor.BondTwoSingletonBaseModel.%
        singletonVerticalCoisometry_reconstruction}
    \leanok
    \uses{thm:mpdo_bond_two_singleton_normalized_algebra}
    The unique label and copy identify the vertical sector coordinate with
    $I$. The corresponding permutation matrix $W$ satisfies
    $WW^\dagger=\Id$. For every $Z\in\MN{2}$, conjugation by $W$ and
    reconstruction by its adjoint give
    \begin{align}
        WZW^\dagger           & =\reindex(Z),
        \notag                                \\
        W^\dagger\reindex(Z)W & =Z.
        \notag
    \end{align}
\end{theorem}
\begin{proof}\leanok
    Each row and column of $W$ contains one nonzero entry, equal to one.
\end{proof}

\begin{theorem}[One-element vertical BNT algebra for the physical tensor]%
    \label{thm:mpdo_bond_two_singleton_tensor_clause}
    \lean{MPOTensor.BondTwoSingletonBaseModel.baseMPOBNTAlgebraTensorClause}
    \lean{MPOTensor.BondTwoSingletonBaseModel.baseMPO_hasBNTAlgebraTensorClause}
    \leanok
    \uses{thm:mpdo_bond_two_singleton_vertical_identity, thm:mpdo_bond_two_singleton_normalized_algebra, thm:mpdo_bond_two_singleton_coisometry}
    The physical tensor $M$ has a vertical BNT with one label, one copy,
    and bond dimension two. Its normalized representative is
    $M_0=\widehat A$, its multiplicity matrix is $\mu_0=(s^{-1})$, and
    $\chi_{0,0,0}=(s)$ gives the structure coefficient
    $c_{0,0,0}^{(L)}=\tr(\chi_{0,0,0}^L)=s^L$ at every positive length.
    The coisometry $W$ identifies the vertical sector coordinates with the
    original physical coordinates. Both the decomposition and its
    reconstruction hold exactly for every vertical letter.
\end{theorem}
